\chapter*{Abstract}
\addcontentsline{toc}{chapter}{Abstract}

Multi-agent language-model systems are often justified by an intuitive division-of-labour argument: several agents can explore more candidate answers than one agent, after which communication and aggregation should select the best result. Final-answer accuracy alone cannot test this mechanism. It conflates failure to discover the answer with failure to retain, endorse, or select an answer that was already present. This dissertation therefore evaluates multi-agent reasoning as a sequence of state transitions from answer discovery to final endorsement.

The study analyses Qwen3-4B in a frozen form and after supervised fine-tuning on MuSiQue and 2WikiMultiHopQA. A single-agent baseline is compared with five multi-agent architectures under matched thinking-token budgets of 1,000 and 2,000 tokens. The primary benchmark contains 9,600 item-configurations, including 8,000 multi-agent configurations and 26,873 intermediate stages. Lexical accuracy is supplemented by a stage-level semantic screen and a preregistered-style blind human audit focused on the two clearest temporal architectures, Debate and Sequential. The audit covers 160 sampled items, 400 non-repeat stage records, and 20 hidden repeats. It separates whether a stage contains sufficient support for the reference answer (Discovery) from whether the stage commits to that answer (Endorsement).

The principal result is that exploration and decision quality separate. Multi-agent systems often expose a correct answer in at least one intermediate output, but later revision, communication, and aggregation can erase that advantage. In the human audit, the weighted harmful transition rate from a correct source endorsement to a wrong final endorsement was 5.45\% for Debate and 0.59\% for Sequential, a difference of 4.86 percentage points. Conversely, beneficial wrong-to-correct transitions were 0.00\% for Debate and 5.82\% for Sequential. A calibrated Mistral judge agreed with human stage labels at 82.5\% for Discovery and 80.5\% for Endorsement, but targeted transition reliability was too weak for it to serve as ground truth. Machine labels are consequently used only for full-trace screening, while the human audit supports the central transition claims.

Increasing the nominal thinking budget also failed to improve aggregate semantic outcomes: discovery fell from 84.19\% to 82.84\%, final accuracy fell from 74.08\% to 72.82\%, and mean wall time increased by 4.83 seconds per item. These results show that additional computation is not equivalent to useful reasoning when architectures repeatedly regenerate, revise, or compress information. The dissertation contributes an auditable evaluation framework, architecture-specific failure localisation, and evidence that multi-agent reasoning should be assessed by both candidate generation and answer preservation rather than final accuracy alone.
